\section{Задание}

Используя хвостовую рекурсию, разработать эффективную программу, позволяющую:
\begin{itemize}
    \item найти длину списка (по верхнему уровню);
    \item найти сумму элементов числового списка;
    \item найти сумму элементов числового списка, стоящих 
    на нечетных позициях исходного списка (нумерация от 0);
    \item сформировать список из элементов числового списка, больших заданного значения;
    \item удалить заданный элемент из списка (один или все вхождения).
    \item объединить два списка.
\end{itemize}

Убедиться в правильности результатов. Для одного из вариантов вопроса 
уметь составить таблицу, отражающую конкретный порядок работы системы.